% =========================================================
% Lattice Theory Breadcrumb
% =========================================================

\subsection*{Why Lattice Theory Gets Its Own Home}

\begin{remark*}[Why lattice theory gets its own home]
Lattice theory begins with a simple observation: in many ordered structures,
two elements have a natural least upper bound and a natural greatest lower
bound. Once those operations exist everywhere, order starts behaving like
algebra.
\end{remark*}

\begin{remark*}[Where the first pass already lives]
The first introduction to lattices is already in
\hyperref[sec:lattices]{the lattice section of the order notes}. That is the
right place for the first encounter, because it grows directly out of partial
orders, suprema, and infima.
\end{remark*}

\begin{remark*}[Why return later]
Once lattices stop being examples and start becoming the main character, the
story gets richer. A fuller lattice-theory chapter will naturally gather
topics such as:
\begin{itemize}
  \item distributive, modular, and Boolean lattices;
  \item complete lattices and closure operators;
  \item sublattices, ideals, and filters;
  \item Galois connections and adjoint situations;
  \item fixed-point theorems such as Knaster--Tarski;
  \item the lattice viewpoint behind topologies, sigma-algebras, and domain theory.
\end{itemize}
\end{remark*}

\begin{remark*}[Why this matters for the rest of the curriculum]
This chapter is a breadcrumb because lattice order quietly reappears all over
the place. Power sets form Boolean lattices, sigma-algebras behave like
closure-stable lattice structures, open sets form a join-rich order, and
fixed-point methods show up in logic, measure, computation, and analysis.
Giving lattice order its own future home makes those later echoes easier to
recognize.
\end{remark*}
